\documentclass[11pt]{article}

\usepackage[utf8]{inputenc}
\usepackage[T1]{fontenc}
\usepackage{lmodern}
\usepackage{geometry}
\usepackage{amsmath,amssymb}
\usepackage{booktabs}
\usepackage{array}
\usepackage{float}
\usepackage{microtype}
\usepackage[hidelinks]{hyperref}

\geometry{margin=1in}

\title{How large is the US data-centre boom?\\
\large Macro and macro-financial ballparks}
\author{}
\date{September 2026}

\begin{document}
\maketitle

\noindent
Three questions for the Moody's roundtable, with a first pass at the arithmetic.
Investment through 2027 is Brandsaas, Garcia, Kurtzman, Nichols and Zytek
(FEDS 2025-109), Dodge vintage September 2025. Census structures through July
2026 are the live check. Electricity is LBNL's 2025 Update. Financing numbers
are BIS, BoE, FSB and the 2025 IG calendar. Anything after 2027, and the 2035
electricity share, is a mechanical hold of a 2026--27 rate.

\section{Questions}

$I_t^{DC}$ is US data-centre GFCF (structures plus installed IT). $I_t$ is a US
investment aggregate. $E_t^{DC}$ is data-centre electricity consumption; $E_t$
is US end-use electricity. $F_t$ is the funds this investment absorbs: retained
earnings spent out of cash and bills, plus bonds, private credit, bank loans
and ABS/CMBS.

\begin{enumerate}
\item Share of investment. What is $s_t^I \equiv I_t^{DC}/I_t$ in 2030 and
2035, against 2023? Report both gross private domestic investment (GPDI) and
private non-residential fixed investment (PNFI).

\item Share of electricity, and the price. What is
$s_t^E \equiv E_t^{DC}/E_t$ in 2030 and 2035, against 2023? Holding other
demand fixed,
\begin{equation}
\frac{\Delta P}{P}
\simeq
\frac{1}{\varepsilon_S}\,
\frac{E_T^{DC}-E_{2023}^{DC}}{E_{2023}}.
\label{eq:price}
\end{equation}
Section~\ref{sec:elas} puts a range on $\varepsilon_S$ instead of setting it
to 1.

\item Funds-market absorption. How large is $F_t$ next to IG credit, private
credit and bank C\&I, and next to the saving pool that prices $r^*$? Can that
move interest rates, and in which market?
\end{enumerate}

\section{Assumptions}

\begin{table}[H]
\centering
\caption{Working assumptions}
\label{tab:assumptions}
\begin{tabular}{>{\raggedright\arraybackslash}p{0.42\textwidth}
>{\raggedright\arraybackslash}p{0.50\textwidth}}
\toprule
Object & Assumption \\
\midrule
$I^{DC}$ through 2027 & Fed Table~2, Q4 annual rate, Dodge vintage
September 2025. 30\% structures, 70\% equipment; 80\% of equipment imported,
so 44\% of $I^{DC}$ is domestic. \\
Live structures check & Census C30, July 2026: \$75.2\,bn SAAR
(+57\% y/y). \\
US investment, 2025 & GPDI \$5.46\,tn (BEA Table 5.2.5). PNFI \$4.25\,tn
(FRED PNFIA). Grown at 2\% real. \\
2030 and 2035 $I^{DC}$ & Hold a 2026--27 Fed rate. Low: 2027 0.25$\times$
inflow (\$357\,bn). Central: 2026 1$\times$ (\$464\,bn). High: 2027
1$\times$ (\$612\,bn). \\
Electricity, 2023--24 & 176\,TWh (4.4\% of 4{,}011\,TWh) and 192\,TWh
(4.7\% of 4{,}110). \\
Electricity, 2030 & LBNL 2025 Update: 649\,TWh, 11.8\% of the NERC 2025
LTRA total (range 9.5--15.3\%; 521--843\,TWh). Data centres are 33\% of
2024--30 load growth. \\
Electricity, 2035 & Hold the 2030 share, or keep the 2023--30 drift to
about 17\%. \\
$\varepsilon_S$ & 0.3 short-run / local peak; 1.0 medium-run national
generation; 2.5 long-run renewables. \\
External-finance share & 30--50\% of $I^{DC}$. \\
Funds-supply elasticity & $\eta=1$; 0.5 if saving is less elastic. \\
\bottomrule
\end{tabular}
\end{table}

\section{The Fed nowcast is a year old}

FEDS 2025-109 has three objects, and only the first two are pinned by data.

\begin{itemize}
\item History, matched to a NIPA imputation (Table~2, Q4 annual rate):
\$10\,bn in 2023, \$58\,bn in 2024, \$179\,bn in 2025.
\item Nowcast from the existing Dodge plans, no assumed future inflows:
mean \$370\,bn by 2026Q2, 5th--95th \$240--450\,bn. Most of the width is
which mega-projects get abandoned (historical $\lambda=0.33$).
\item Forecast through 2027, which needs an inflow assumption: \$357\,bn
if new plans run at one-quarter of the 2024--25 pace, \$612\,bn if they
continue, \$921\,bn if they double. After about a year the current planning
stock no longer sets the level.
\end{itemize}

The Dodge extract is September 2025. We are now in the year where the nowcast
runs out and the inflow assumption takes over. Census structures kept rising:
\$50.7\,bn SAAR in April 2026, \$70.8\,bn in June, \$75.2\,bn in July
(+57\% y/y). At the paper's 30/70 split, \$75\,bn of structures is about
\$250\,bn all-in, under the \$370\,bn 2026Q2 mean. Equipment is more
backloaded than the shell, so the gap can still close; the September 2025
mean can also just have been high. Dodge Momentum Index commercial planning
fell in June 2026 on a dip in data-centre filings, and Wood Mackenzie pipeline
additions have fallen for three quarters from a 2025Q3 peak. Use
\$370--460\,bn as a 2026 working range.

The 2$\times$ path to \$921\,bn is dropped below. By 2027 the paper is
simulating new plans, not reading the current stock.

\section{Investment intensity}

From the 2025 BEA levels, at 2\% real,
\begin{equation}
\begin{aligned}
I_{2030}^{GPDI} &= 5.46\times 1.02^{5}=\$6.03\,\text{tn}, &
I_{2035}^{GPDI} &= \$6.65\,\text{tn},\\
I_{2030}^{PNFI} &= 4.25\times 1.02^{5}=\$4.69\,\text{tn}, &
I_{2035}^{PNFI} &= \$5.18\,\text{tn}.
\end{aligned}
\end{equation}

\begin{table}[H]
\centering
\small
\caption{Data-centre share of US investment}
\label{tab:inv}
\begin{tabular}{lrrrr}
\toprule
 & $I^{DC}$ (\$bn) & \% GPDI & \% PNFI 2030 & \% PNFI 2035 \\
\midrule
2023 Q4 (Fed) & 10 & 0.2 & --- & --- \\
2025 Q4 (Fed) & 179 & 3.3 & 4.2 (2025) & --- \\
2026, 1$\times$ inflow & 464 & 8.5 of 2025 & --- & --- \\
Census $\times$ 10/3, Jul 2026 & $\approx$ 250 & 4.6 & 5.9 & --- \\
\addlinespace
Low (hold 2027 0.25$\times$) & 357 & 5.9 / 5.4 & 7.6 & 6.9 \\
Central (hold 2026 1$\times$) & 464 & 7.7 / 7.0 & 9.9 & 9.0 \\
High (hold 2027 1$\times$) & 612 & 10.2 / 9.2 & 13.0 & 11.8 \\
\bottomrule
\end{tabular}
\end{table}

The 2023 and 2025 GPDI shares use \$5.46\,tn. In the low / central / high
rows the GPDI pair is 2030 then 2035.

The central hold is about 8\% of GPDI and 10\% of PNFI in 2030, from 0.2\%
of GPDI in 2023. Pre-boom is 2023, not the 2024Q4 Fed print of \$60\,bn.
Spain is 2.5--3.5\% of GFCF in the companion note, so the US intensity is
three to four times larger.

Eight percent of GPDI is about 1.4\% of GDP if $I/Y\approx 0.18$. Only 44\%
of $I^{DC}$ is domestic, so the GDP impulse is closer to 0.6\% of GDP. The
Fed's own 2026 contribution, import-adjusted, is 0.33--0.49 percentage points
of nominal GDP growth (Table~2). Equipment depreciates fast, so a 10\% PNFI
flow share is not a 10\% capital-stock share.

\section{Electricity intensity}
\label{sec:elec}

LBNL 2025 Update: 649\,TWh in 2030 is 11.8\% of the NERC 2025 LTRA forecast,
not of 2023 EIA end use. That implies a 2030 total of about 5{,}500\,TWh.
Non-data-centre load rises 926\,TWh (24\%) from 2024 to 2030. Data centres
are 33\% of load growth over those six years.

\begin{table}[H]
\centering
\caption{Electricity share}
\label{tab:elecshare}
\begin{tabular}{lrrr}
\toprule
 & $E^{DC}$ (TWh) & Share of $E$ & $\Delta E^{DC}/E_{2023}$ \\
\midrule
2023 & 176 & 4.4\% & --- \\
2024 & 192 & 4.7\% & $+$0.4\% \\
2030, LBNL low & 521 & 9.5\% & $+$8.6\% \\
2030, LBNL central & 649 & 11.8\% & $+$11.8\% \\
2030, LBNL high & 843 & 15.3\% & $+$16.6\% \\
2035, hold 2030 share & $\approx$ 680 & 11.8\% & $\approx$ $+$12\% \\
2035, share-drift high & $\approx$ 980 & $\approx$ 17\% & $\approx$ $+$20\% \\
\bottomrule
\end{tabular}
\end{table}

In the central 2030 case the data-centre share goes from 4.4\% to 11.8\%.
That ratio does not need an elasticity.

\section{Supply elasticity}
\label{sec:elas}

Electricity supply is layered, and the layers differ.

With $\varepsilon_S=1$, equation~\eqref{eq:price} maps a 12\% demand shift
into a 12\% national price. That describes a medium-run, national, annual
energy market: other demand frozen, new generation available at rising cost.

There is no clean US estimate of the elasticity of delivered kWh at
data-centre nodes. By layer:

\begin{itemize}
\item Short-run operating and peak. Wholesale demand elasticities sit near
0--0.35. Existing wind's operating elasticity is about 0.1. At the peak the
supply curve is the merit order plus a vertical capacity constraint; PJM
capacity prices going from \$29/MW-day to the cap are that market. Use
$\varepsilon_S\approx 0.3$ or below.

\item Medium-run generation, by fuel. BOEM MarketSim elasticities, from EIA
AEO price cases: 1.50 (gas), 0.27 (coal), 0.53 (nuclear), 2.03 (solar),
0.65 (onshore wind), 0.68 (other non-hydro renewables). Weighted by the 2024
US mix (gas 43\%, nuclear 18\%, coal 15\%, wind 10.5\%, utility solar 5\%),
that is about 0.9--1.1. Unit-elastic is a fair medium-run national guess.

\item Long-run renewable capacity. A Georgia Tech IV using state RPS
schedules puts the price elasticity of renewable capacity near 2.7. That is
capacity entry, not kWh at a congested node. It is the right order of
magnitude once generation can adjust and transmission is not the bind.

\item This shock, in existing models. IMF WP/25/81 (April 2025) has US
electricity prices $+$0.9\% under current policies and $+$8.6\% when
renewable expansion and transmission are constrained. If the quantity shock
is LBNL-central 12\%, the implied elasticities are about 13 and 1.4.
CMU/NC~State's $+$8\% on US generation cost by 2030 is the constrained case
again.
\end{itemize}

The bind for data centres is interconnection, transformers (two to three
years) and substations (four to six years), not US generating-fuel supply.
The IMF constrained case is the national benchmark; local peak markets can
print a multiple of that.

\begin{table}[H]
\centering
\caption{National price map, LBNL-central 12\% quantity shift}
\label{tab:elas}
\begin{tabular}{lrrr}
\toprule
Layer & $\varepsilon_S$ & $\Delta P/P$ & Reading \\
\midrule
Short-run / local peak & 0.3 & $+$39\% & Merit order plus capacity; PJM-like \\
Medium-run national generation & 1.0 & $+$12\% & AEO-weighted fuel mix \\
IMF constrained T\&D & $\approx$ 1.4 & $+$8.6\% & Model output \\
Long-run renewables & 2.5 & $+$5\% & Capacity entry if the wires exist \\
IMF current policies & $\approx$ 13 & $+$0.9\% & Generation and T\&D keep up \\
\bottomrule
\end{tabular}
\end{table}

For a national number, use about 10\% by 2030 when wires and renewables lag:
the IMF constrained case is $+$8.6\%, the unit-elastic map is $+$12\%. If
generation and wires keep up, about $+$1\%. In the local peak markets the
energy session will discuss, several times the national figure. Keep
$\varepsilon_S=1$ for the national annual market; use 0.3 or below for PJM.

\section{Funds-market absorption}

Split the \$460\,bn of investment into retained earnings spent and funds raised.
\begin{equation}
I_t^{DC}
=
\underbrace{RE_t}_{\text{retained earnings, cash and bills run down}}
+
\underbrace{F_t}_{\text{bonds, private credit, bank loans, ABS/CMBS}}.
\end{equation}
$r^*$ and crowding-out respond to $\Delta I^{DC}$. Credit spreads respond to
$F$. From Fed 2023Q4 (\$10\,bn) to the central 2026 hold (\$464\,bn),
\begin{equation}
\Delta I^{DC}\approx \$450\,\text{bn per year}.
\end{equation}
The BoE FSR has private credit at 34\% of AI-investment financing in 2025,
from 9\% in 2024. Hyperscalers issued about \$121\,bn of US bonds in 2025
(four times the 2020--24 average). Data-centre debt issuance was about
\$182\,bn. At a 30--50\% external share,
\begin{equation}
F \approx \$140\text{--}230\,\text{bn per year}
\end{equation}
at the central rate. Vontobel's \$300\,bn of AI and data-centre bonds over
the next year, and \$1.5\,tn over five years, is the bond slice only.
Barclays forecasts \$945\,bn of net US corporate issuance in 2026, so
\$300\,bn is again about a third of the net calendar.

\begin{table}[H]
\centering
\caption{Absorbed funds against the markets that price them}
\label{tab:funds}
\begin{tabular}{>{\raggedright\arraybackslash}p{0.30\textwidth}
>{\raggedright\arraybackslash}p{0.28\textwidth}
>{\raggedright\arraybackslash}p{0.32\textwidth}}
\toprule
Market $M_j$ & Comparison & Size \\
\midrule
US GPDI / saving, \$5.46\,tn & $\Delta I^{DC}=\$450$\,bn & 8\% of the US
saving flow \\
Global saving, order of \$25\,tn & $\Delta I^{DC}=\$450$\,bn & about 2\%
of global saving \\
Net USD IG issuance & AI-related share, 2025 & about 30\% of the net
calendar; about \$300\,bn expected next year \\
Hyperscaler + DC debt & \$121\,bn and \$182\,bn in 2025 & large
single-sector bid \\
Private credit stock & $>$\$200\,bn AI-related of about \$2.5\,tn
(BIS) & about 8\% of the stock; more than a third of 2025 deals (FSB).
BIS projects \$300--600\,bn by 2030 \\
Bank C\&I to data centres & \$14.9\,bn through 2025Q3 (MSCI) & negligible \\
AI-adjacent bank C\&I & \$450\,bn commitments, \$150\,bn drawn & 0.8\% of
bank assets (Chicago Fed) \\
Broader AI-company debt & about \$1.2\,tn (JPM); \$243\,bn of AI-firm
bonds in 2025 vs \$3.0\,tn NFC issuance (BIS/OECD) & large for a sector;
small against all corporate debt \\
DC ABS/CMBS outstanding & about \$50--60\,bn & well measured; small \\
\bottomrule
\end{tabular}
\end{table}

On the FSB's \$1.5--2.0\,tn denominator this would be 10--13\% of private
credit. BIS later puts outstanding private credit near \$2.5\,tn, and
AI-related loans at 8\% of that stock. More than a third of 2025
originations is the tighter comparison.

\begin{equation}
\frac{\Delta r}{r}
\simeq
\frac{1}{\eta}\,\frac{\Delta I^{DC}}{S}.
\label{eq:rstar}
\end{equation}
At $\eta=1$ and $S=\$5.46$\,tn, \$450\,bn of extra investment demand raises
$r$ by 8\% of itself. At a 4\% real rate, about 30 basis points. At
$\eta=0.5$, about 65 basis points. Against global saving of order \$25\,tn,
about 7 basis points at $\eta=1$.

Those are $r^*$ increments, inside the usual $r^*$ standard error. Put
10--30\,bp in a natural-rate discussion. Treasury yields have a much larger
denominator.

Thirty percent of net dollar IG issuance is already a thick calendar.
Private credit already prices this (8\% of the stock, a third of 2025
deals). Banks barely hold it.

Spending retained earnings still uses national saving: cash spent on GPUs
is cash not held in bills. What changes is who holds the risk. BIS
Bulletin 120 documents capex outgrowing free cash flow. Circular
lab--hyperscaler--chipmaker financing inflates gross issuance relative to
net absorption, so do not add the \$1.2\,tn, \$200\,bn and \$182\,bn
headlines.

\section{Working answers}

\begin{table}[H]
\centering
\caption{Ballparks}
\label{tab:answers}
\begin{tabular}{>{\raggedright\arraybackslash}p{0.34\textwidth}
>{\raggedright\arraybackslash}p{0.58\textwidth}}
\toprule
Question & Number \\
\midrule
Share of US investment, 2030 / 2035 & Central: about 8\% / 7\% of
GPDI, 10\% / 9\% of PNFI. 2023: 0.2\% of GPDI. \\
Share of US electricity vs 2023 & 4.4\% to 11.8\% in 2030
(9.5--15.3). One-third of 2024--30 load growth. 2035: about 12\% or
about 17\%. \\
National electricity price & $+$9 to $+$12\% by 2030 if wires and
renewables lag. About $+$1\% if they keep up. Local peak: a multiple. \\
Funds taken out & $\Delta I^{DC}\approx\$450$\,bn/yr;
external $F\approx\$140$--230\,bn/yr. \\
$r^*$? & About 7--30\,bp against global / US saving at $\eta=1$. \\
Credit rates? & About 30\% of net USD IG; about 8\% of the
private-credit stock, more than a third of 2025 deals. Banks: no. \\
\bottomrule
\end{tabular}
\end{table}

An adverse scenario that matches these magnitudes is constrained electricity
plus private-credit and SPV refinancing.

\section{Limitations}

The 2030 and 2035 investment numerators hold a 2026--27 Fed rate; they are
not in FEDS 2025-109. Wood Mackenzie's 195\,GW of signed construction or
supply agreements (26\% of 2025 US peak load) is a capacity upper bound,
not GFCF or TWh. A 30--50\% haircut still belongs on utility-stated demand;
LBNL is already a haircut relative to those queues.

The price map is partial equilibrium. It leaves out demand destruction,
large-load tariffs, behind-the-meter gas (about 2\,GW online against 50\,GW
of IT load), and the split between energy and capacity prices. The
generation-weighted elasticity of 0.9--1.1 averages fuels. It is not the
elasticity of power at a congested Dominion or ERCOT node.

The funds comparison freezes 2025 financing shares. If hyperscaler free
cash flow recovers, $F$ shrinks and the credit-market pressure eases; the
$r^*$ arithmetic does not.

\section*{Sources}

Brandsaas, Garcia, Kurtzman, Nichols and Zytek, ``Estimating Aggregate Data
Center Investment with Project-level Data,'' FEDS 2025-109 (Dodge vintage
September 2025; Table~2).

Census C30, July 2026 release (1 September 2026): private data-centre VIP
\$75.166\,bn SAAR.

BEA / FRED: GPDI \$5.459\,tn in 2025 (Table 5.2.5); PNFIA \$4.251\,tn.

Smith, Hubbard, Newkirk, Ganeshalingam, Holecek, Sartor, Mills and Shehabi,
\emph{United States Data Center Energy Usage Report: 2025 Update}, LBNL,
June 2026. Denominator: NERC 2025 LTRA.

Bogmans, Gomez-Gonzalez, Ganpurev, Melina, Pescatori and Thube, ``Power Hungry:
How AI Will Drive Energy Demand,'' IMF WP/25/81.

BOEM MarketSim own-price supply elasticities (AEO-derived, September 2021
update). Georgia Tech IV estimate of the long-run renewable-capacity elasticity
(about 2.7).

BIS Bulletin 120; BIS Bulletin 130; BIS ``Financing the digital economy''
(2026); BoE FSR (July 2026); FSB private-credit report (May 2026); Vontobel and
Barclays on the 2025--26 IG calendar; Chicago Fed (February 2026).

\end{document}
